% !TeX root = ../main.tex
\section{Correctness Validation}\label{sec:correctness}

Every build variant is checked two independent ways before any of its timings are trusted
(\code{make test}, \code{scripts/validate.sh}).

\emph{Reference comparison.} \code{src/reference.c} contains plain scalar loops compiled with
\code{-fno-tree-vectorize -ffp-contract=off}, so it evaluates each kernel in strict, non-fused,
left-to-right order. The driver runs the kernel under test and the reference on identical inputs and
compares element-by-element (or the scalar result, for the dot product) before timing.

\begin{knobs}
  Cases  & kernels $\times$ \{float32, float64\}; $N$ around multiples of the lane count
           ($1,2,3,7,8,9,15,16,17,31,33,64,127,1000,1024,4095,4096,4097,65536$); start offsets
           $4\ldots\qty{60}{\byte}$; strides $2,4,8,16$; gather; int32 multiply \\
  Tolerance & relative error $\le 10^{-5}$ (float32) and $\le 10^{-12}$ (float64). The bound is loose
             enough for the one legal source of disagreement---FMA contraction in the vectorized
             builds versus the non-fused reference---and tight enough to catch real bugs. \\
  Command & \code{make test} \\
\end{knobs}

\emph{Known-answer self-test} (\code{simd_bench --selftest}), independent of the reference: the dot
product of two all-ones vectors equals $N$; $y=2x+y$ with $x=1,y=3$ gives $5$ everywhere; $z=x\,y$
with $x=2,y=3$ gives $6$; and the stencil with $a=b=c=1$ on a constant-1 input gives $3$ in the
interior. These guard against a reference and kernel sharing the same mistake.

The sweep covers sizes with and without a vector-width remainder (to exercise the scalar/masked tail)
and deliberately misaligned start addresses (to exercise the loop prologue), so correctness is
established precisely in the regimes the later experiments stress.

Result: all \num{570} validation cases pass for every variant (scalar, autovec, native); the maximum
observed float32 relative error was below $10^{-6}$, consistent with FMA contraction alone.
